\documentclass[11pt,a4paper]{article}

\usepackage[T1]{fontenc}
\usepackage[utf8]{inputenc}
\usepackage{lmodern}
\usepackage{microtype}
\usepackage{amsmath,amssymb,amsthm,mathtools}
\usepackage{geometry}
\usepackage{xcolor}
\usepackage{hyperref}
\usepackage{enumitem}
\usepackage{booktabs}
\usepackage{titlesec}

\geometry{margin=1in}
\setlength{\parindent}{0pt}
\setlength{\parskip}{0.62em}
\linespread{1.035}
\hypersetup{colorlinks=true,linkcolor=blue!45!black,citecolor=blue!45!black,urlcolor=blue!45!black}
\titleformat{\section}{\Large\bfseries\color{blue!35!black}}{\thesection.}{0.7em}{}
\titleformat{\subsection}{\large\bfseries\color{black}}{\thesubsection.}{0.6em}{}
\setlist[itemize]{leftmargin=1.35em,itemsep=0.2em,topsep=0.25em}
\setlist[enumerate]{leftmargin=1.45em,itemsep=0.2em,topsep=0.25em}

\newtheoremstyle{soft}%
  {0.7em}{0.7em}{\itshape}{}%
  {\bfseries\color{blue!35!black}}{.}{0.5em}{}
\theoremstyle{soft}
\newtheorem{theorem}{Theorem}[section]
\newtheorem{lemma}[theorem]{Lemma}
\newtheorem{proposition}[theorem]{Proposition}
\newtheorem{conjecture}[theorem]{Conjecture}

\theoremstyle{definition}
\newtheorem{definition}[theorem]{Definition}
\newtheorem{remark}[theorem]{Remark}

\newcommand{\C}{\mathbb C}
\newcommand{\CP}{\mathbb{CP}}
\newcommand{\Q}{\mathbf Q}
\newcommand{\z}{\mathbf z}
\newcommand{\w}{\mathbf w}
\newcommand{\norm}[1]{\left\lVert #1\right\rVert}
\newcommand{\partline}[2]{\section*{#1\\[-0.3em]{\large #2}}\addcontentsline{toc}{section}{#1: #2}}
\newcommand{\idea}[1]{\begin{center}\fcolorbox{blue!20!black}{blue!4}{\begin{minipage}{0.90\linewidth}\small #1\end{minipage}}\end{center}}

\title{Universitas Pandrosion\\\large A Short Unified Series on Smale, Quotients, and Geometric Root Finding}
\author{Ivan Besevic}
\date{April 2026 -- compact reading edition}

\begin{document}
\maketitle

\begin{abstract}
This text rewrites Parts I--IX of the Pandrosion--Smale notes as a single short manuscript.  The guiding object is the Pandrosion quotient
\[
   Q_P(a,b)=\frac{P(b)-P(a)}{b-a},
\]
and, for polynomial systems, its matrix version
\[
   F(b)-F(a)=\Q_F(a,b)(b-a).
\]
The same quotient gives three connected stories: a non-local root-finding operator, a compact language for Smale's Mean Value Conjecture, and an algorithmic route toward Smale's 17th problem based on starts, scaling, and line-search geometry.  The emphasis here is readability.  Proofs are given when the series has clean identities; conditional or empirical claims are stated as such.
\end{abstract}

\tableofcontents

\section*{Reading principle}
\addcontentsline{toc}{section}{Reading principle}
Pandrosion is best read as a change of viewpoint.  Newton differentiates at one point.  Pandrosion compares two points and asks for the slope that makes the comparison exact.  In one variable that slope is a quotient.  In several variables it is a matrix.  The series below develops what happens when this exact slope is used both as an algorithmic step and as a geometric language for Smale-type questions.

\idea{The central rule is not to linearise first, but to compare first.  The comparison
\(F(b)-F(a)=\Q_F(a,b)(b-a)\) is exact for polynomials.  All later constructions are ways of choosing the two points \(a,b\), scaling the geometry, or accelerating the resulting motion.}

\partline{Part I}{The Pandrosion operator}

Let \(P\in\C[z]\).  For \(a\neq b\), define
\begin{equation}\label{eq:scalar-Q}
   Q_P(a,b)=\frac{P(b)-P(a)}{b-a}.
\end{equation}
The value \(Q_P(a,b)\) is the exact slope of the secant from \(a\) to \(b\).  If \(b\to a\), it becomes \(P'(a)\).  Thus Newton's method is the infinitesimal edge of a larger non-local picture.

For systems \(F:\C^n\to\C^n\), the scalar quotient is replaced by a slope matrix.

\begin{definition}[Pandrosion slope matrix]
A matrix \(\Q_F(a,b)\in\C^{n\times n}\) is a Pandrosion slope matrix between \(a,b\in\C^n\) if
\begin{equation}\label{eq:matrix-Q}
   F(b)-F(a)=\Q_F(a,b)(b-a).
\end{equation}
For polynomial systems one canonical choice is obtained by expanding each monomial telescopically in the coordinates.
\end{definition}

Once \(\Q_F(a,b)\) is available, the fixed-anchor Pandrosion step is
\begin{equation}\label{eq:pandrosion-step}
   \mathcal P_{F,a}(b)=a-\Q_F(a,b)^{-1}F(a),
\end{equation}
whenever the matrix is invertible.  If \(a\) is near a regular zero and \(b\) is chosen in a compatible local geometry, this step moves toward the zero.  If \(b\to a\), then \(\Q_F(a,b)\to DF(a)\), and \eqref{eq:pandrosion-step} becomes Newton's step.

The original root-extraction formula appears as a special case.  For \(x\neq0\), \(p\ge2\), and
\[
   S_p(s)=1+s+\cdots+s^{p-1},
\]
the map
\begin{equation}\label{eq:root-map}
   h(s)=1-\frac{x-1}{xS_p(s)}
\end{equation}
has fixed points \(s=x^{-1/p}\) up to the \(p\) complex branches.  This is the one-dimensional Pandrosion construction written in the language of a geometric sum.

\begin{proposition}[Local picture]
Let \(\zeta\) be a regular zero of \(F\).  In a sufficiently small neighbourhood of \(\zeta\), the slope matrix \(\Q_F(a,b)\) is close to \(DF(\zeta)\).  Hence the Pandrosion step is locally well defined and fixes \(\zeta\).
\end{proposition}

The series repeatedly uses one additional idea: scale the geometry before applying the quotient.  Instead of solving \(F(z)=0\), solve
\begin{equation}\label{eq:scaled-system}
   G(y)=F(Ay)=0,
\end{equation}
where \(A\) is a homothetic or linear change of coordinates.  The zero is then recovered by \(z=Ay\).  This is not cosmetic.  It changes the shape of the basins in which the quotient acts.

\partline{Part II}{Smale's Mean Value Conjecture as a quotient bound}

Smale's Mean Value Conjecture concerns a polynomial \(P\) of degree \(d\), a point \(z\) with \(P'(z)\neq0\), and the critical points \(c\) of \(P\).  In Pandrosion notation the conjectural inequality becomes
\begin{equation}\label{eq:smale-mvc}
   \min_{P'(c)=0}\left|\frac{Q_P(z,c)}{P'(z)}\right|
   \le \frac{d-1}{d}.
\end{equation}
Since \(P'(z)=Q_P(z,z)\), the conjecture says: among the critical points, at least one gives a secant quotient no larger than \((d-1)/d\) of the infinitesimal quotient at \(z\).

\idea{In this form, a critical point is not only a point where \(P'\) vanishes.  It is a possible Pandrosion target.  Smale's conjecture asks whether one critical target is always good enough.}

For quadratics the statement is exact.  If \(P(z)=(z-\alpha)(z-\beta)\), then at the unique critical point \(c=(\alpha+\beta)/2\),
\[
   \frac{Q_P(z,c)}{P'(z)}=\frac12.
\]
For cubics, the sharp constant is \(2/3\), attained by the symmetric extremal family in the limiting configuration.  More generally, the product over all critical points is controlled by a resultant.  If \(c_1,\ldots,c_{d-1}\) are the roots of \(P'\) and \(R(w)=Q_P(z,w)\), then
\begin{equation}\label{eq:critical-product}
   \prod_{j=1}^{d-1}Q_P(z,c_j)=\frac1d\operatorname{Res}(R,P').
\end{equation}
This identity is the first sign that the quotient formulation is not merely notation: it packages all critical targets into one algebraic object.

\partline{Part III}{The vanishing identity}

Normalize the conjecture at the origin: assume
\[
   P(0)=0,\qquad P'(0)=1.
\]
Write the roots as \(0,\alpha_2,\ldots,\alpha_d\), and set
\[
   R(z)=\prod_{k=2}^{d}(z-\alpha_k).
\]
The Lagrange--Sylvester identity gives
\begin{equation}\label{eq:vanishing-identity}
   1=-\sum_{k=2}^{d}\frac{1}{\alpha_k R'(\alpha_k)}.
\end{equation}
This is a root-side constraint.  It does not prove the Mean Value Conjecture by itself, but it explains why an extremal configuration must be highly balanced.

At a critical point \(c\), one has
\begin{equation}\label{eq:critical-R}
   R(c)=-cR'(c).
\end{equation}
Thus the Smale ratio can be read through the value of \(R\) or through the derivative of the root polynomial at \(c\).  The symmetric family
\begin{equation}\label{eq:zd-z}
   P(z)=z^d-z
\end{equation}
attains the value \((d-1)/d\) exactly.  This family becomes the natural model extremal in the following parts.

\partline{Part IV}{Fibers and the Pandrosion inverse}

At a critical point \(c\), the value \(P(c)\) has a fiber: all points \(z\) satisfying \(P(z)=P(c)\).  Removing the critical point itself, call the remaining points the co-fiber of \(c\).  The quotient from \(c\) back to the origin is
\begin{equation}\label{eq:inverse-quotient}
   Q_P(c,0)=\frac{P(c)}{c}.
\end{equation}
Thus the Smale ratio is an inverse Pandrosion quotient.

If the co-fiber points are \(z_1,\ldots,z_{d-2}\), then the product decomposition is
\begin{equation}\label{eq:cofiber-product}
   \left|\frac{P(c)}{c}\right|=|c|\prod_{j=1}^{d-2}|z_j|.
\end{equation}
This turns a value estimate into a geometric estimate on a fiber.

Two related identities complete the picture.  Let \(\alpha_0,\ldots,\alpha_{d-1}\) be the roots of \(P\).  At a critical point \(c\),
\begin{align}
   \sum_{j=0}^{d-1}\frac{1}{(\alpha_j-c)P'(\alpha_j)}&=-\frac{1}{P(c)},\label{eq:fiber-bridge}\\
   \sum_{j=0}^{d-1}\frac{1}{(\alpha_j-c)^2P'(\alpha_j)}&=0.\label{eq:fiber-vanish}
\end{align}
The first equation bridges the fiber to the value \(P(c)\).  The second is an orthogonality relation.  Together they suggest that the conjecture is controlled by the geometry of fibers, not only by the location of critical points.

For normalized polynomials
\[
   P(z)=z^d+a_{d-1}z^{d-1}+\cdots+a_1z,
\]
the Pandrosion reduction gives, at every critical point,
\begin{equation}\label{eq:pandrosion-reduction}
   \frac{P(c)}{c}=\frac{\widetilde q(c)}{d},\qquad
   \widetilde q(z)=\sum_{k=1}^{d-1}k a_{d-k}z^{d-1-k}.
\end{equation}
Smale's bound becomes the search for a critical point with
\[
   |\widetilde q(c)|\le d-1.
\]

\partline{Part V}{The quartic slice}

For degree four, write
\begin{equation}\label{eq:quartic}
   P(z)=z^4+az^3+bz^2+z.
\end{equation}
Then the reduction \eqref{eq:pandrosion-reduction} becomes
\begin{equation}\label{eq:quartic-q}
   \frac{P(c)}{c}=\frac{\widetilde q(c)}4,
   \qquad \widetilde q(z)=az^2+2bz+3.
\end{equation}
At critical points, the same quantity simplifies to
\begin{equation}\label{eq:quartic-simplify}
   \widetilde q(c)=2\bigl(1-c^2(a+2c)\bigr).
\end{equation}
Hence the quartic case of the conjecture is equivalent to the following compact statement: for every \((a,b)\in\C^2\), the cubic equation
\[
   4c^3+3ac^2+2bc+1=0
\]
has a root \(c\) satisfying
\begin{equation}\label{eq:quartic-target}
   |1-c^2(a+2c)|\le\frac32.
\end{equation}

The symmetric extremal \(P(z)=z^4-z\) gives equality.  The product of the three values \(\widetilde q(c_k)\) is explicitly
\begin{equation}\label{eq:quartic-resultant}
   \prod_{k=1}^{3}\widetilde q(c_k)=4a^3-a^2b^2-18ab+4b^3+27.
\end{equation}
The part isolates a clean cubic problem.  The remaining obstruction is to show that this cubic in the critical values must always enter the disk of radius \(3/2\).

\partline{Part VI}{The quintic slice}

For the centered quintic slice
\begin{equation}\label{eq:quintic}
   P(z)=z^5+az^3+bz^2+z,
\end{equation}
the critical equation has four roots \(c_j(a,b)\).  The reduced quotient values satisfy a resultant product formula
\begin{equation}\label{eq:quintic-product}
   \prod_{j=1}^{4}\widetilde q(c_j)=\frac{N(a,b)}{625},
\end{equation}
where
\begin{equation}\label{eq:Nab}
   N(a,b)=16a^4-4a^3b^2-128a^2+144ab^2-27b^4+256.
\end{equation}
At \((a,b)=(0,0)\), i.e. \(P(z)=z^5+z\), all four critical points satisfy
\begin{equation}\label{eq:quintic-extremal}
   |\widetilde q(c_j)|=\frac45.
\end{equation}
This is the expected extremal value.

The part supplies three pieces of evidence around this extremal:
\begin{itemize}
\item a local symbolic certificate that the symmetric point is a candidate local maximum;
\item an asymptotic descent program suggesting that large \((a,b)\) cannot preserve the extremal value;
\item numerical verification on dense parameter grids.
\end{itemize}
The conclusion is deliberately modest: the centered quintic slice is reduced to an explicit two-parameter algebraic problem, with a strong candidate proof strategy but not a closed global proof in every singular direction.

\partline{Part VII}{The algorithmic branch: descent without trusting the basin}

The MVC branch studies critical points.  The algorithmic branch returns to root finding.  The goal is to find approximate zeros with controlled cost, in the spirit of Smale's 17th problem.

The base iteration is a Pandrosion-accelerated orbit.  But an orbit may stagnate before entering a good basin.  Part VII therefore adds a descent safeguard.  Given an anchor \(z_0\), first try the accelerated Pandrosion candidate.  If it fails to reduce the residual, replace it by a derivative-free Armijo-type step based on finite quotients.

A probe radius \(\varepsilon\) gives directional quotients
\begin{equation}\label{eq:probe-quotient}
   Q_{\varepsilon,u}(z_0)=\frac{P(z_0+\varepsilon u)-P(z_0)}{\varepsilon u},
   \qquad u\in\{1,-1,i,-i\}.
\end{equation}
These approximate \(P'(z_0)\) without explicitly differentiating.  With the modulus-gradient identity
\begin{equation}\label{eq:modgrad}
   \frac{d}{dt}\bigg|_{t=0}|P(z_0-tv)|^2
   =-2\operatorname{Re}\!\bigl[v\overline{P(z_0)}P'(z_0)\bigr],
\end{equation}
one chooses a descent direction and backtracks until the residual decreases.

Under the usual regularity assumptions, this fallback gives quantified descent.  Conditional on non-degeneracy of the starts and basin-entry estimates, the resulting multi-start scheme yields polynomial BSS-type bounds in the univariate case and suggests an analogous multivariate route through the matrix quotient \(\Q_F\).  These statements are algorithmic companions to the known solutions of Smale's 17th problem, not replacements for them.

\partline{Part VIII}{Starts, spirals, and the role of \texorpdfstring{$T_2$}{T2}}

The next question is not only how to move, but where to start.  The Cauchy circle contains the roots, but equispaced starts can resonate with the root configuration.  Part VIII replaces the rigid ring by a geometric spiral:
\begin{equation}\label{eq:strategy-B}
   z_k=\rho q^k e^{ik\varphi},
   \qquad q\in(0,1),\qquad \varphi=\pi(3-\sqrt5).
\end{equation}
The radius decreases geometrically, while the golden angle spreads the phases with low discrepancy.

The empirical message is simple.  Among the tested accelerated maps \(T_n\), the \(T_2\) level with frequent reanchoring is the best practical compromise.  It is cheap enough to repeat, yet strong enough to enter basins rapidly on the tested families.  Combined with Strategy B, it suppresses the observed phase-transition failures that affected the earlier Cauchy-ring strategy.

This part is partly empirical and partly programmatic.  Its role is to identify the start geometry that makes the quotient iteration reliable before the local contraction theory takes over.

\partline{Part IX}{Extended line search and the undershooting gap}

The final part revisits a failure mode.  On Kostlan--Smale inputs, a Newton-like step from a typical point near the unit scale has magnitude
\begin{equation}\label{eq:ks-newton-mag}
   |P(z)/P'(z)|\asymp d^{-1/2},
\end{equation}
while nearby roots live at distance \(\Theta(1)\).  Backtracking can only shrink a step that is already too small.  This creates the undershooting gap.

The fix is to allow forwardtracking.  Instead of testing only \(\tau\le1\), define
\begin{equation}\label{eq:ELS}
   \mathcal T_{\mathrm{ELS}}=\{2^k:k_-\le k\le k_+\},
\end{equation}
and choose
\begin{equation}\label{eq:ELS-step}
   z_{n+1}=z_n-\tau^\star\Delta_n,
   \qquad
   \tau^\star=\arg\min_{\tau\in\mathcal T_{\mathrm{ELS}}}|P(z_n-\tau\Delta_n)|,
   \qquad \Delta_n=\frac{P(z_n)}{P'(z_n)}.
\end{equation}
The optimal \(\tau\) on KS is predicted to scale like \(\sqrt d\), matching the gap in \eqref{eq:ks-newton-mag}.  Numerically this changes pre-basin motion from a long crawl into a short jump.

Conditional on the stated probabilistic decorrelation and optimal-\(\tau\) estimates, the Pandrosion \(B'\)-\(T_2\)-ELS scheme has expected cost
\begin{equation}\label{eq:dlogd}
   \mathbb E[\mathrm{cost}]=O(d\log d)
\end{equation}
on the univariate Kostlan--Smale ensemble.  The contribution is the mechanism: a symmetric line search repairs the geometric undershooting left by backtracking-only descent.

\section*{Closing view}
\addcontentsline{toc}{section}{Closing view}
The nine parts are one path through the same idea.  The quotient \(Q\) first appears as a secant.  It becomes a root-finding operator, then a language for critical points, then an inverse-fiber geometry, then an algorithmic tool.  Scaling, starts, and line searches do not replace the quotient; they make the quotient act in a better geometry.

The proved identities are algebraic: telescoping quotients, product formulae, vanishing identities, and local reductions.  The global MVC program remains partly open beyond the low-degree slices and verified reductions.  The algorithmic claims likewise separate definitions, conditional bounds, and empirical certificates.  The unifying claim of the series is more structural: Pandrosion gives a single geometric vocabulary in which Smale's critical-point inequality and Smale-style root-finding algorithms can be written naturally.

\begin{thebibliography}{9}
\bibitem{smale1981} S. Smale, \emph{The fundamental theorem of algebra and complexity theory}, Bull. Amer. Math. Soc. 4 (1981), 1--36.
\bibitem{smale1998} S. Smale, \emph{Mathematical problems for the next century}, Math. Intelligencer 20 (1998), 7--15.
\bibitem{shubsmale} M. Shub and S. Smale, \emph{Complexity of Bezout's theorem I: Geometric aspects}, J. Amer. Math. Soc. 6 (1993), 459--501.
\bibitem{beltranpardo} C. Beltran and L. M. Pardo, \emph{Smale's 17th problem: average polynomial time to compute affine and projective solutions}, J. Amer. Math. Soc. 22 (2009), 363--385.
\bibitem{lairez} P. Lairez, \emph{A deterministic algorithm to compute approximate roots of polynomial systems in polynomial average time}, Found. Comput. Math. 17 (2017), 1265--1292.
\bibitem{edelmankostlan} A. Edelman and E. Kostlan, \emph{How many zeros of a random polynomial are real?}, Bull. Amer. Math. Soc. 32 (1995), 1--37.
\bibitem{schmidt} W. M. Schmidt, \emph{Diophantine approximation}, Lecture Notes in Mathematics 785, Springer, 1980.
\end{thebibliography}

\end{document}
